% ============================================
% DOC-22 (ZH)
% 中文翻译 — SPEC-935-R356
% ============================================
\section*{\textbf{DOC-22} --- 内部备忘录第47/2026号}
\markboth{DOC-22 --- 内部备忘录第47/2026号}{}

\noindent \textit{从鉴定档案中提取的文件。}

\vspace{0.5em}\noindent\rule{\textwidth}{0.5pt}\vspace{0.5em}


\noindent {\footnotesize\texttt{公共档案馆 $\cdot$ 书证 $\cdot$ 案卷1,260}}\\[-0.2em]

\noindent \textbf{米纳斯吉拉斯州公共档案馆 —— 文献管理局}\\
\noindent \textbf{致:} 特藏部。 \textbf{抄送:} 行政总署。\\
\noindent \textbf{事由:} 结束对登记号1,260物品的分析。\\
\noindent\textbf{日期:} 2026年8月1日。

\noindent 各位：

\noindent 1. 鉴于登记号1,260（巴尔巴塞纳收容医院）名下笔记本物品已交由法医精神病学家卢西亚·门德斯医生（CRM-MG 28,391）进行外部分析，且该专业人员已在意见书完成前申请临床休假；

\noindent 2. 鉴于重新诊断报告（DOC-08）已附入原始病历，且目前没有任何司法、行政或学术决定要求该物品继续流通；

\noindent 3. 鉴于本机构的文献安全准则与馆藏稳定原则；

\noindent \textbf{兹决定:}

\noindent a) 立即终止对1,260物品的公开查阅；

\noindent b) 将该物品存放于永久保险柜，不设新的访问安排；

\noindent c) 通知所有相关方，本案已归档，不再受理新的调阅申请。

\noindent \textbf{理由:} 该物品不构成犯罪物证，不属于任何进行中的司法程序，也不具备足以证明让技术团队承担风险的例外历史价值。其持续流通损害文献安全与相关专业人员的心理完整性。

\noindent \textbf{附加建议:} 建议管理局评估不将此备忘录录入文档管理信息系统的适当性。信息化的记录会使文件可以被检索。而可检索的文件，可能被再次申请调阅。

\noindent \textit{罗纳尔多·托雷斯}\\
\noindent \textit{米纳斯吉拉斯州公共档案馆 —— 文献管理局局长}
\noindent \textit{页首手写批注，蓝色墨水，草书:}
\noindent \textquotedblleft{\textit{罗纳尔多——未能办到。备忘录发出的次日早晨，物品已从保险柜中被取走。追踪系统记录到凌晨3点14分保险柜被打开。监控中没有识别出任何人。门是密封的。封条完好。按您的要求归档结案。”}}

\noindent “—— 特藏部，2026年8月2日。”





\vspace{0.5em}
\noindent\textit{\textbf{↪ 路线:} \rota{DOC-01} \rotasep \rota{LUC-02} \rotasep \rota{DOC-23}}
